\DocumentMetadata{
  pdfversion=2.0,
  pdfstandard=ua-2,
  lang=en-US,
  tagging=on,
  tagging-setup={math/setup=mathml-SE,tabsorder=structure}
}
\documentclass[11pt,letterpaper]{article}
\usepackage[margin=0.68in,includefoot]{geometry}
\usepackage{newcomputermodern}
\usepackage{amsmath,amscd,mathtools}
\providecommand{\square}{\mdlgwhtsquare}
\usepackage[hidelinks]{hyperref}
\hypersetup{
  pdftitle={Topology PhD Exam, August 2012},
  pdfauthor={University of Florida Department of Mathematics},
  pdfsubject={Graduate mathematics examination},
  pdfdisplaydoctitle=true,
  bookmarksopen=true
}
\setcounter{secnumdepth}{0}
\setlength{\parindent}{0pt}
\setlength{\parskip}{0.28em}
\setlength{\emergencystretch}{3em}
\setlength{\leftmargini}{2.15em}
\setlength{\leftmarginii}{2.65em}
\setlength{\leftmarginiii}{3em}
\renewcommand{\labelenumi}{\textbf{\theenumi.}}
\pagestyle{empty}
\raggedbottom
\allowdisplaybreaks
\newenvironment{examproblems}[1]{%
  \begin{enumerate}%
  \setcounter{enumi}{#1}%
  \setlength{\topsep}{0.35em}%
  \setlength{\partopsep}{0pt}%
  \setlength{\itemsep}{0.48em}%
  \setlength{\parsep}{0.18em}%
}{\end{enumerate}}
\newenvironment{examsubparts}{%
  \begin{enumerate}%
  \setlength{\topsep}{0.18em}%
  \setlength{\partopsep}{0pt}%
  \setlength{\itemsep}{0.16em}%
  \setlength{\parsep}{0.1em}%
}{\end{enumerate}}
\newenvironment{examinstructions}{%
  \par\begingroup\itshape
}{%
  \par\endgroup\vspace{0.18em}
}
\makeatletter
\renewcommand\section{\@startsection{section}{1}{0pt}%
  {-1.25ex plus -.25ex minus -.1ex}%
  {0.55ex plus .1ex}%
  {\normalfont\large\bfseries}}
\renewcommand{\@maketitle}{%
  \newpage\null\vskip -1.2em%
  {\footnotesize\bfseries
    \href{https://www.ufl.edu/}{University of Florida}, \href{https://math.ufl.edu/}{Department of Mathematics}\par}%
  \vskip 0.18em%
  \tagstructbegin{tag=Title}%
    {\LARGE\bfseries\href{https://gma.math.ufl.edu/exams/topology-phd/}{Topology PhD Exam}, August 2012\par}%
  \tagstructend%
  \vskip 0.35em%
  \tagmcbegin{artifact}%
    \hrule height 1pt%
  \tagmcend%
  \vskip 0.55em%
}
\makeatother
\title{Topology PhD Exam, August 2012}
\author{}
\date{}
\begin{document}
\maketitle
\thispagestyle{empty}
\begin{examinstructions}
Work the following problems and show all work. Support all statements to the best of your ability.
\end{examinstructions}
\begin{examproblems}{0}
\item
Prove that the 3-sphere \(S^3\) is not homeomorphic to the 3-space \(\mathbb{R}^3\).
\item
Let \(A \subset \mathbb{R}^2\) be an infinite countable subspace.
\begin{examsubparts}
\item[\textnormal{(a)}]
Can~\(A\) be connected?
\item[\textnormal{(b)}]
Is \(\mathbb{R}^2 \setminus A\) connected?
\end{examsubparts}
\item
Construct a map \(f:T^2 \to S^2\) of degree 3 where \(T^2=S^1 \times S^1\) is a torus.
\item
Show that there is no map \(S^2 \to T^2\) of degree 3.
\item
Let \(GL_n\) be the space of all invertible real \(n \times n\)-matrices. Is \(GL_n\) compact? Is it connected?
\end{examproblems}
\begin{examinstructions}
Answer the following with complete definitions or statements or short proofs.
\end{examinstructions}
\begin{examproblems}{5}
\item
Is the 2-sphere \(S^2\) with~\(k\) points removed homeomorphic to a topological group for
\begin{examsubparts}
\item[\textnormal{(a)}]
\(k=1\)?
\item[\textnormal{(b)}]
\(k=2\)?
\item[\textnormal{(c)}]
\(k=3\)?
\item[\textnormal{(d)}]
\(k=0\)?
\end{examsubparts}
\item
State the Lefschetz Fixed Point Theorem.
\item
Compute the Euler characteristic \(\chi(\mathbb{CP}^5 \times \mathbb{RP}^5 \times S^5 \times S^2)\).
\item
State the Baire Category Theorem
\item
Does every function \(f:\mathbb{N} \to \mathbb{R}\) admit a continuous extension \(\bar{f}:\beta\mathbb{N} \to \mathbb{R}\) to the Stone–Čech compactification?
\item
State the Five Lemma.
\item
What can you say about the~\(k\)-th cohomology group of a closed oriented manifold for
\begin{examsubparts}
\item[\textnormal{(a)}]
\(k=n\)?
\item[\textnormal{(b)}]
\(k=n-1\)?
\end{examsubparts}
\item
Does there exist a covering space of the figure eight that has a non-trivial abelian fundamental group?
\item
Describe all connected subsets of the real line~\(\mathbb{R}\).
\item
State the Contraction Mapping Theorem.
\end{examproblems}
\end{document}
